\chapter{Additive Number Theory and Sieves}\label{additive-number-theory}
\chapterstatus{planned}

Additive number theory asks which integers are sums of elements of a given set --- squares, $k$-th
powers, primes --- and in how many ways. Its two great tools are the circle method, which turns a
counting problem into an integral over the circle, and sieve methods, which count integers free of
small prime factors. The chapter ends with additive combinatorics, where the sets are arbitrary.

\noindent\textbf{Prerequisites.} Chapter~\ref{number-theory} (Elementary Number Theory), Chapter~\ref{prime-numbers} (Prime Numbers), Chapter~\ref{analytic-number-theory} (Analytic Number Theory), Chapter~\ref{harmonic-analysis} (Harmonic Analysis) and Chapter~\ref{combinatorics} (Combinatorics).

\noindent\textbf{Planned contents.}
\begin{enumerate}
\item Sums of squares: the numbers $r_2(n)$ and $r_4(n)$ (Jacobi), Gauss's circle problem and its error term.
\item Waring's problem: Hilbert's theorem via Linnik's or Hilbert's method; the numbers $g(k)$ and $G(k)$.
\item The Hardy--Littlewood circle method: major and minor arcs, Weyl sums and Weyl's inequality, the
singular series and the singular integral; Waring's problem for large $s$.
\item Vinogradov's three primes theorem; the binary Goldbach problem and what is known.
\item Elementary sieves: the sieve of Eratosthenes--Legendre, Brun's sieve and Brun's theorem on twin
primes, the Selberg sieve, the large sieve inequality.
\item The Bombieri--Vinogradov theorem (statement), Chen's theorem (statement), bounded gaps between
primes (Zhang, Maynard; survey).
\item Schnirelmann density and Schnirelmann's theorem; Mann's theorem.
\item Additive combinatorics: the Cauchy--Davenport theorem, sumsets, the Pl\"unnecke--Ruzsa
inequalities, Freiman's theorem (statement), Roth's theorem on three-term progressions,
Szemer\'edi's and the Green--Tao theorems (survey).
\item Partitions: the Hardy--Ramanujan asymptotic formula (via the circle method) and Ramanujan's congruences.
\item Exercises with solutions.
\end{enumerate}
